\scanpage{62}
faisceau
\[
\mathcal F^{(\nu)}:\quad U\longmapsto F(U^{(\nu+1)}/U)=\varprojlim_i F(\Delta_U^{(\nu)}(i))
\]
sur $X_{\mathrm{zar}}$, soit $\mathcal F^{(\nu)}$, et que pour $\nu$ variable, ces faisceaux forment un faisceau cosimplicial sur $X_{\mathrm{zar}}$. Ceci posé, on aura un isom
\[
H^*(X_{\mathrm{strat}},F)\simeq\mathbb H^*(X_{\mathrm{zar}},\underset{\substack{\Vert\\(\nu\mapsto\mathcal F^{(\nu)})}}{\mathcal F^*}),
\]
qui implique une suite spectrale d'hypercoh
\[
\boxed{H^*(X_{\unsure{\mathrm{inf}|\mathrm{strat}}},F)\Longleftarrow
E_2^{pq}=H^p(\deleted{\unsure{F}}\nu\longmapsto H^q(X_{\mathrm{zar}},\mathcal F^{(\nu)}))}
\]
\unsure{généralisant} l'isomorphisme donné plus haut.

\deleted{\unsure{Regardons} le complexe $\mathcal F^*$, prenons pour simplifier $F=\mathcal O_{X_{\mathrm{strat}}}$.}

\margin{\uline{Coh. cristalline}}
6.b \unsure{Lorsqu'on travaille} en cohomologie cristalline, on ne peut (si $X$ n'est pas \unsure{f\textsuperscript{t}} lisse sur $S$) affirmer que $\widetilde X$ couvre l'objet final. Mais supposons que l'on puisse trouver une immersion
\mspage{61}
